% ===================================================================
%  5. TAULA — Etxea eta nola eraikitzen den.
%  Traduction 2026 d'apres texto/io, controlee sur texto/fr, texto/en
%  et texto/es. Consigne : voir texto/eu/10-taula-01.tex.
%
%  Le tableau 5 est un DIALOGUE, et les deux volumes composent le nom
%  du parleur en petites capitales : on garde \textsc{}.
%
%  Les renvois du plan sont des LETTRES — (a) le couloir, (b) le
%  salon... — et non des numeros. On les ecrit comme l'ido les ecrit,
%  y compris la ou l'imprimeur a compose « (1) » pour « (l) ».
%
%  LE TABLEAU DES METIERS, ET LE BASQUE LES NOMME TOUS PAR LE MEME
%  PROCEDE : la chose faite, plus « -gin » ou « -zain ». Le macon est
%  « hargina », l'homme de la pierre ; le charpentier « zurgina »,
%  l'homme du bois ; le forgeron « errementaria », l'homme du fer ; le
%  vitrier « beirategilea », celui qui fait le verre. Une seule regle
%  de formation la ou les langues indo-europeennes en ont dix.
% ===================================================================

%%K t05-apar-1 apar
\VUsaut{6.0mm}
\VUtitre{15.0pt}{60}{5. TAULA}
\VUsaut{1.4mm}
\VUtitre{11.4pt}{0}{Etxea eta nola eraikitzen den.}
\VUsaut{2.2mm}

%%K t05-tit-1 sub
\VUpk{10.2pt}{40}{Topaketa zoriontsua.}
\VUsaut{3.0mm}

%%K t05-01-1 p
\textsc{Joanes}. --- Osaba, oso gogo handia dut jakiteko nola eraikitzen den
\VUgras{etxe} bat.

%%K t05-01-2 p
\textsc{Osaba} \textsuperscript{(80)}. --- Ez da zaila, mutiko; zatoz nirekin eta
erakutsiko dizut Bernard jaunak \textsuperscript{(79)} eraikitzen ari dena eta
zeinetan biziko bainaiz.

%%K t05-01-3 p
\textsc{Joanes}. --- Eta nork marraztu du etxe honen \VUgras{plano}
\textsuperscript{(76)}?

%%K t05-01-4 p
\textsc{Osaba}. --- \VUgras{Arkitektoak} \textsuperscript{(75)}; marrazki oso
argia egin zuen, onartu egin zioten, eta \VUgras{kontratistak}
\textsuperscript{(77)} lanak hasi zituen.

%%K t05-01-5 p
\textsc{Joanes}. --- Eta nor da kontratista?

%%K t05-01-6 p
\textsc{Osaba}. --- Blanc jauna, zure lagun Jakoberen aita. Berak zuzentzen ditu
langileak Roux jaunarekin, arkitektoarekin, batera.

%%K t05-01-7 p
Hara hor bera! Blanc jauna garaiz dator bere \VUgras{erregelarekin}
\textsuperscript{(78)}, eta Roux jauna bere planoarekin.

%%K t05-01-8 p
Egun on, jaunak. Zer moduz zaudete?

%%K t05-01-9 p
\textsc{Roux jauna}. --- Ederki, eskerrik asko. Egun on, Joanes; zeinen hazi den
mutikoa!

%%K t05-01-10 p
\textsc{Osaba}. --- Egia da, gizontxo bat da jada. Oso serioa da; ikusten duen
guztia azaldu behar zaio. Gaur etxe bat nola eraikitzen den jakin nahi du.

%%K t05-tit-2 sub
\VUpk{10.2pt}{40}{Langileak.}
\VUsaut{3.0mm}

%%K t05-01-11 p
\textsc{Blanc jauna}. --- Zatoz nirekin, bada, adiskide gazte, eta oraintxe
azalduko dizut dena, oso maite baititut ikasi nahi duten haurrak.

%%K t05-01-12 p
Ikusten duzu langile hori, \VUgras{pala} \textsuperscript{(82)} eskuan duena:
\VUgras{lubakigilea} \textsuperscript{(81)} da. \VUgras{Erretena}
\textsuperscript{(46)} zulatzen ari da, ur-hodia jartzeko. Berak hondeatu zuen
lurra etxea dagoen tokian, harginak lau \VUgras{horma} altxa zitzan. Horma horien
lurpeko zatiari \VUgras{zimendua} \textsuperscript{(2)} deitzen zaio. Zimenduak
ixten duen espazioak \VUgras{soto} bat \textsuperscript{(3)} osatzen du. Soto
honek leihotxo bat du, \VUgras{arnasleiho} \textsuperscript{(5)} deitua,
\VUgras{ate} bat \textsuperscript{(4)} eta eskailera bat.

%%K t05-01-13 p
\VUgras{Harginak} \textsuperscript{(108)} hormak altxatzen jarraitu zuen,
\VUgras{ateentzat} \textsuperscript{(8)} eta \VUgras{leihoentzat}
\textsuperscript{(17)} hutsuneak utziz.

%%K t05-01-14 p
\textsc{Joanes}. --- Baina ez dut ikusten, hargin hori!

%%K t05-01-15 p
\textsc{Blanc jauna}. --- Etxean lanean amaitu du jada. Hantxe dago,
\VUgras{ukuiluaren} \textsuperscript{(54)} ondoan, bere \VUgras{aldamioan}
\textsuperscript{(110)}, \VUgras{garbitegiaren} \textsuperscript{(56)} horma
altxatzen.

%%K t05-01-16 p
Paleta bat, aska bat eta \VUgras{plomua} \textsuperscript{(109)} ditu. Baina izen
horiek ez dituzu gogoan hartuko.

%%K t05-01-17 p
\textsc{Joanes}. --- Bai horixe, gogoan hartuko ditut, osabari paleta txiki bat
eta aska bat eskatuko baitizkiot, eta plomua neuk egingo dut soka batekin.

%%K t05-tit-3 sub
\VUsaut{3.0mm}
\VUpk{10.2pt}{40}{Materialak.}
\VUsaut{3.0mm}

%%K t05-01-18 p
\textsc{Osaba}. --- Ondo da. Baina esan ezazu, Joanes, zer behar da etxe bat
eraikitzeko?

%%K t05-01-19 p
\textsc{Joanes}. --- \VUgras{Harri landua} \textsuperscript{(59)} eta
\VUgras{morteroa} \textsuperscript{(65)} behar dira.

%%K t05-01-20 p
\textsc{Osaba}. --- Hala da. Begira ezazu kapela handiko gizon hori, bere
\VUgras{gurdian} \textsuperscript{(136)}. \VUgras{Gurdizaina}
\textsuperscript{(135)} da. Egunero \VUgras{harri-koskorrak}
\textsuperscript{(60)} ekartzen ditu hona. Gurdia patioan hustutzen du, eta
\VUgras{harginek} \textsuperscript{(83)} koskorrak lantzen dituzte eta beren
\VUgras{harri-zerraz} \textsuperscript{(85)} mozten. \VUgras{Laguntzaileak}
\textsuperscript{(146)} harginari eramaten dizkio, eta hark bere lekuan jartzen
ditu.

%%K t05-01-21 p
Bitartean \VUgras{morterogileak} \textsuperscript{(87)} morteroa prestatzen du.
\VUgras{Harea} \textsuperscript{(62)}, \VUgras{karea} \textsuperscript{(63)} eta
\VUgras{ura} \textsuperscript{(64)} behar ditu. Karea haren ondoan dago
\VUgras{zaku} batean \textsuperscript{(71)}; \VUgras{harginaren morroiak}
\textsuperscript{(111)} harea ekartzen dio \VUgras{eskorgan}
\textsuperscript{(112)}, eta \VUgras{ikastuna} \textsuperscript{(113)} ponpan
\textsuperscript{(58)} dago. \VUgras{Ontzia} \textsuperscript{(114)} betetzen du.
Morteroa laster egongo da prest. \VUgras{Emailea} \textsuperscript{(107)} hartu
eta eskaileretatik aldamiora eramaten du, non harginak etxearen alde hau
amaitzen ari baita.

%%K t05-01-22 p
Eta orain: nola deitzen da \VUgras{teilatua} \textsuperscript{(31)} egiten duen
langilea?

%%K t05-01-23 p
\textsc{Joanes}. --- \VUgras{Teilatugilea} \textsuperscript{(118)} da.

%%K t05-tit-4 sub
\VUsaut{3.0mm}
\VUpk{10.2pt}{40}{Etxeko teilatua.}
\VUsaut{3.0mm}

%%K t05-01-24 p
\textsc{Blanc jauna}. --- Bai, baina lehenik \VUgras{zurgina}
\textsuperscript{(115)} dator.

%%K t05-01-25 p
Zurginak \VUgras{gapirioak} \textsuperscript{(35)} egiten ditu. \VUgras{Habeak}
\textsuperscript{(37)} lotzen ditu \VUgras{ziriz} \textsuperscript{(117)}. Goiko
langile horiek, belusezko galtza zabaletan, zurginak dira. Batek zur-puska
eusten du, besteak ziria lantzen du bere \VUgras{aizkoraz}
\textsuperscript{(116)}. \VUgras{Tximiniaren} \textsuperscript{(41)} ondokoa
teilatugilea da. Latak \VUgras{habeetara} (*) iltzatzen ditu, eta
\VUgras{arbela} \textsuperscript{(66)} lateetara. Batzuetan teila jartzen du.

%%K t05-noto-1 noto
\VUnotes{143.2mm}{%
(*) \VUgras{Balk-o} = alem. \textit{Balken}, ingel. \textit{joist}, fran.
\textit{solive}, ital. \textit{travicello}, errus. \textit{балка}, esp. eta
port. \textit{viga}. Erro teknikoa, oraindik onartu gabea, baina hemen
proposatzen dena frantsesezko \textit{solive} hitzaren esanahia emateko, zeina
ez baita ez \textit{trabo} (fran. \textit{poutre}) ez zur-puska. Zoru eta sabaia
eusten dituen enbor landua da.%
}

%%K t05-01-26 p
Ukuiluko teilatua \VUgras{teilazkoa} \textsuperscript{(55)} da, etxekoa
\VUgras{arbelezkoa} \textsuperscript{(32)}, eta berandarena
\VUgras{zinkezkoa} \textsuperscript{(24)}.

%%K t05-01-27 p
\textsc{Joanes}. --- Zinkezko teilatuak ere teilatugileak egiten ditu?

%%K t05-01-28 p
\textsc{Blanc jauna}. --- Ez, \VUgras{latorrigilea} \textsuperscript{(121)} edo
\VUgras{iturgina} da, hor etxeko teilatuan \VUgras{argizuloa}
\textsuperscript{(38)} jartzen ari dena. Berak jartzen ditu \VUgras{erretenak}
\textsuperscript{(43)} eta \VUgras{ur-hodiak} \textsuperscript{(42)} ere. Zinka
eta latorria soldatzen ditu. Berak finkatu zituen \VUgras{tximistorratza}
\textsuperscript{(40)} eta \VUgras{haize-orratza} \textsuperscript{(39)}
\VUgras{frontoian} \textsuperscript{(36)}.

%%K t05-01-29 p
\textsc{Joanes}. --- Beraz, etxea prest dago?

%%K t05-tit-5 sub
\VUsaut{3.0mm}
\VUpk{10.2pt}{40}{Tresnak.}
\VUsaut{3.0mm}

%%K t05-01-30 p
\textsc{Blanc jauna}. --- Ez guztiz. \VUgras{Igeltseroak}
\textsuperscript{(99)} \VUgras{adreiluz} \textsuperscript{(61)} egiten ditu
trenkadak, zeinek etxea gela desberdinetan banatzen baitute. \VUgras{Igeltsua}
\textsuperscript{(70)} jartzen du \VUgras{sabaietan} \textsuperscript{(26)} eta
hormetan. Orain \VUgras{berandaren} \textsuperscript{(33)} sabaian ari da lanean.
Harginak bezala, \VUgras{paleta} \textsuperscript{(100)} eta \VUgras{aska}
\textsuperscript{(101)} ditu.

%%K t05-01-31 p
Gero arotzak ateak, leihoak, \VUgras{zurezko zoruak} \textsuperscript{(16)} eta
\VUgras{leiho-kontrak} \textsuperscript{(19)} egiten ditu.

%%K t05-01-32 p
Orain patioan ari da lanean bere \VUgras{lan-mahaiaren} \textsuperscript{(90)}
aurrean. \VUgras{Zerra} bat \textsuperscript{(94)} du, \VUgras{zura}
\textsuperscript{(69)} mozteko, \VUgras{arraspa} \textsuperscript{(91)},
\VUgras{estutzailea} \textsuperscript{(92)}, \VUgras{zizela} \textsuperscript{(94
\textit{bis})}, \VUgras{konpasa} \textsuperscript{(95)} eta zurezko
\VUgras{mailua} \textsuperscript{(95 \textit{bis})}.

%%K t05-01-33 p
\textsc{Joanes}. --- A, arotza izatea harginarena baino ere alaiagoa da! Osaba,
lan-mahai bat, zerra bat eta arraspa bat erosiko dizkidazu? Txakurrarentzat
etxola bat egingo dut.

%%K t05-01-34 p
\textsc{Osaba}. --- Erosiko dizkizut, mutiko.

%%K t05-01-35 p
\textsc{Joanes}. --- Zein pozik nagoen! Eta nor da sarrerako atean dagoen gizon
hori?

%%K t05-tit-6 sub
\VUsaut{3.0mm}
\VUpk{10.2pt}{40}{Margo-lanak.}
\VUsaut{3.0mm}

%%K t05-01-36 p
\textsc{Blanc jauna}. --- \VUgras{Margolaria} \textsuperscript{(102)} da. Bere
eskaileran dago, \VUgras{margo} ontzi batekin \textsuperscript{(103)} eta bere
\VUgras{pintzelarekin} \textsuperscript{(104)}. Sarrerako atearen zura estaltzen
du, luzaroago iraun dezan.

%%K t05-01-37 p
Berak paperatzen ditu gelak ere.

%%K t05-01-38 p
\VUgras{Paperjartzaileak} \textsuperscript{(107)} \VUgras{eskailera txikian}
\textsuperscript{(106)}, \VUgras{balkoian} \textsuperscript{(28)},
\VUgras{gortina} \textsuperscript{(29)} zintzilikatzen du.

%%K t05-01-39 p
Leiho handiaren ondoan dagoen langilea \VUgras{beirategilea}
\textsuperscript{(96)} da. \VUgras{Kristalak} \textsuperscript{(18)} jartzen ditu
\VUgras{mastikaz} \textsuperscript{(73)}. Eta sototako atean belauniko dagoen
langile hori \VUgras{sarrailagilea} \textsuperscript{(97)} da. \VUgras{Sarraila}
\textsuperscript{(6)} jartzen ari da. \VUgras{Lima} \textsuperscript{(98)} haren
ondoan datza.

%%K t05-01-40 p
\textsc{Joanes}. --- Eta burdin puska bat \VUgras{mailuz}
\textsuperscript{(84)} jotzen ari den hori ere sarrailagilea da?

%%K t05-01-41 p
\textsc{Blanc jauna}. --- Zehatzago esanda, \VUgras{errementaria}
\textsuperscript{(125)} da. \VUgras{Burdinazko piezak} \textsuperscript{(67)}
forjatzen ditu \VUgras{elektrizistarentzat} \textsuperscript{(124)}, zeina
\VUgras{kotxetegiaren} \textsuperscript{(51)} teilatuan baitago.
\VUgras{Sutegia} \textsuperscript{(126)}, \VUgras{ingudea}
\textsuperscript{(127)} eta \VUgras{kurrikak} \textsuperscript{(128)} ditu.

%%K t05-01-42 p
Elektrizistak telefonoaren \VUgras{kobrezko haria} \textsuperscript{(44)}
finkatzen du.

%%K t05-tit-7 sub
\VUpk{10.2pt}{40}{Lorategiko lana.}
\VUsaut{3.0mm}

%%K t05-01-43 p
\textsc{Osaba}. --- \VUgras{Patioaren} \textsuperscript{(45)} barrenean,
\VUgras{burdinsarearen} \textsuperscript{(47)} eta \VUgras{atalasearen}
\textsuperscript{(48)} ondoan, \VUgras{lorezaina} \textsuperscript{(129)} ikusten
duzu. \VUgras{Lorategi} txiki bat \textsuperscript{(49)} prestatzen ari da. Lurra
irauli du bere \VUgras{aitzurraz} \textsuperscript{(131)}, hazia ereiten du eta
\VUgras{eskuarez} \textsuperscript{(130)} berdintzen. Gero bere
\VUgras{ureztontzitik} \textsuperscript{(132)} \VUgras{lorategi-lerroak}
\textsuperscript{(150)} ureztatuko ditu, eta landareak haziko dira.

%%K t05-01-44 p
\VUgras{Zaldizainak} \textsuperscript{(133)}, zaldia garbitu eta bazka eman
berri dionak, \VUgras{kotxea} \textsuperscript{(52)} garbitzen du belakiaz,
\VUgras{esku-ponpaz} \textsuperscript{(134)} eta ur-ontzi batez. Gero kotxetegian
sartuko du eta atea \VUgras{morroilo} astun batez \textsuperscript{(53)} itxiko
du.

%%K t05-01-45 p
Hara, orain pozik zaudela uste dut; azalpen aski izan da.

%%K t05-01-46 p
\textsc{Joanes}. --- Bai, osaba, baina etxea barrutik ikusi nahi nuke.

%%K t05-01-47 p
\textsc{Osaba}. --- Hori bihar izango da. Roux jaunak planoa azalduko dizu eta
gela bakoitza izendatuko.

%%K t05-tit-8 sub
\VUsaut{3.0mm}
\VUpk{10.2pt}{40}{Etxearen banaketa.}
\VUsaut{2.4mm}

%%K t05-apar-2 apar
{\centering\textit{(Ikus planoa.)}\par}
\VUsaut{2.4mm}

%%K t05-01-48 p
\textsc{Arkitektoa}. --- Goazen, Joanes, etxearen zati desberdinetatik eramango
zaitut.

%%K t05-01-49 p
Sar gaitezen \VUgras{beheko solairuan} \textsuperscript{(7)}. Hona hemen
\VUgras{txirrinaren} \textsuperscript{(11)} botoia. Hiru \VUgras{maila}
\textsuperscript{(14)} igotzen ditugu, \VUgras{ataripekoak}
\textsuperscript{(9)}; \VUgras{atalasea} \textsuperscript{(10)} zeharkatzen
dugu, \VUgras{sarrerako atearena} \textsuperscript{(8)} --- eta hona gaude
\VUgras{eskaileraren}
\textsuperscript{(13)} oinean.

%%K t05-01-50 p
\textsc{Joanes}. --- Hau korridorea da.

%%K t05-01-51 p
\textsc{Arkitektoa}. --- Ez, hau \VUgras{atalondoa} \textsuperscript{(12)} da.
\VUgras{Korridorea} \textsuperscript{(a)} barrenean dago. Eskuinean
\VUgras{egongela} \textsuperscript{(b)} eta \VUgras{jangela}
\textsuperscript{(c)}; jangelaren aurrean \VUgras{sukaldea}
\textsuperscript{(d)} eta \VUgras{jaki-gela} \textsuperscript{(e)}; gero,
fatxadaren aldetik, \VUgras{bulegoa} \textsuperscript{(f)}. \VUgras{Beranda}
\textsuperscript{(23)}, bere \VUgras{beirazko atearekin} \textsuperscript{(25)},
egongelari itsatsita dago.

%%K t05-01-52 p
Orain igo gaitezen eskaileretatik. Eutsi \VUgras{eskudelari}
\textsuperscript{(15)} eta konta itzazu mailak. Hamalau dira. Hona hemen
\VUgras{atalburua} \textsuperscript{(m)}: lehen \VUgras{solairuan}
\textsuperscript{(27)} gaude.

%%K t05-tit-9 sub
\VUsaut{3.0mm}
\VUpk{10.2pt}{40}{Lehen solairua.}
\VUsaut{3.0mm}

%%K t05-01-53 p
Eskaileraren aurrean \VUgras{komuna} \textsuperscript{(20)} eta
\VUgras{apaingela} \textsuperscript{(j)} daude. Eskuinean \VUgras{logela}
\textsuperscript{(g)} ederra, etxearen \VUgras{fatxadara}
\textsuperscript{(1)} begira dauden bi leihorekin. Ezkerrean \VUgras{bainugela}
\textsuperscript{(1)} eta \VUgras{gonbidatuen gela} \textsuperscript{(i)},
\VUgras{alkoba} handi batekin \textsuperscript{(h)}. Logelaren ondoan
\VUgras{haurren gela} \textsuperscript{(k)}. \VUgras{Terraza}
\textsuperscript{(21)} zabalera ematen du. Terraza horren azpian beste komun bat
\textsuperscript{(20)}, eta hormaren ondoan \VUgras{ezkila}
\textsuperscript{(22)}, bazkarira deitzeko jotzen dena.

%%K t05-01-54 p
\textsc{Joanes}. --- Igo gaitezen \VUgras{bigarren solairura}.

%%K t05-01-55 p
\textsc{Arkitektoa}. --- Ez dago bigarren solairurik. Teilatupean
\VUgras{mansardak} \textsuperscript{(33)} eta ganbara \textsuperscript{(34)}
baino ez daude. Itzulia amaitu da, jaits gaitezke.

%%K t05-01-56 p
\textsc{Joanes}. --- Eskerrik asko, jauna, oso interesgarria izan da. Aitari
ikusi dudan guztia kontatuko diot.
\VUsaut{4.0mm}
\VUfilet{22mm}
